% ----------------------------------------------------------
% RESUMOS
% ----------------------------------------------------------

% resumo em português
\setlength{\absparsep}{18pt} % ajusta o espaçamento dos parágrafos do resumo
\begin{resumo}
A inteligência artificial agêntica desloca o software do papel de ferramenta que apoia o trabalho humano para o de sistema que executa trabalho de forma autônoma, e a adoção dessa tecnologia avança mais rápido do que as estruturas capazes de torná-la confiável. Este trabalho parte dessa lacuna entre capacidade técnica e confiabilidade e propõe a Arquitetura PASE (\textit{Personalized Adaptive Software Ecosystem}), uma arquitetura híbrida humano-IA que automatiza o ciclo de vida do software com agentes de IA sob supervisão humana. Os requisitos da arquitetura derivam de uma análise multidimensional do ecossistema de software, conduzida por revisão narrativa estruturada de 22 trabalhos sobre indústria, desenvolvimento, experiência do usuário, segurança e governança, e mercado de trabalho, que resultou em cinco requisitos explícitos. A PASE organiza-se em três camadas: a Camada de Tarefas, que capta demandas de usuários, da equipe e da própria IA e as filtra com memória do que já foi implementado e recusado; a Camada de Orquestração, que prioriza, modela dependências e posiciona a decisão humana sobre a tarefa, o plano, a entrega e a publicação; e a Camada de Execução, que produz código em ambiente isolado, verificado por máquina e revisado por um agente independente, e que confina a personalização da interface em tempo de execução a zonas com contratos verificáveis. Para validar a proposta, foi implementada uma prova de conceito, aplicada à evolução de um painel financeiro pessoal. Uma demanda em linguagem natural percorreu as três camadas e foi publicada na versão 1.0.0 do produto, com decisão humana em cada ponto crítico, e o agente de personalização adaptou a interface por usuário e recusou, no pior caso, três conteúdos maliciosos que contornaram a validação do servidor. A validação também expôs falhas reais: metade das entregas aprovadas antes da introdução da revisão independente foi revertida, e as estimativas de custo do planejamento erraram por mais de uma ordem de grandeza. Cada falha foi absorvida pela arquitetura na forma de um mecanismo verificável. Como contribuições, o trabalho oferece a análise multidimensional e seus requisitos, a especificação da arquitetura, a prova de conceito com suas evidências e um princípio de projeto que emergiu da validação: garantias que dependem apenas de instrução ao modelo não são garantias, e devem ser verificadas por código ou decididas por uma pessoa.

\textbf{Palavras-chave}: Inteligência Artificial Agêntica. Engenharia de Software. Autonomia Supervisionada. Sistemas Multiagente. Interfaces Adaptativas. Governança de IA.
\end{resumo}

% resumo em inglês
\begin{resumo}[Abstract]
 \begin{otherlanguage*}{english}
Agentic artificial intelligence is shifting software from a tool that supports human work to a system that performs work autonomously, and its adoption is advancing faster than the structures needed to make it trustworthy. This work addresses that gap between technical capability and trustworthiness and proposes the PASE Architecture (Personalized Adaptive Software Ecosystem), a hybrid human-AI architecture that automates the software life cycle with AI agents under human supervision. The architecture's requirements derive from a multidimensional analysis of the software ecosystem, conducted through a structured narrative review of 22 works on industry, development, user experience, security and governance, and the labor market, which yielded five explicit requirements. PASE is organized into three layers: the Task Layer, which captures demands from users, the team and the AI itself and filters them against a memory of what has already been implemented and rejected; the Orchestration Layer, which prioritizes tasks, models dependencies and places human decisions over the task, the plan, the delivery and the release; and the Execution Layer, which produces code in an isolated environment, verified by machine and reviewed by an independent agent, and which confines runtime interface personalization to zones with verifiable contracts. To validate the proposal, a proof of concept was implemented and applied to the evolution of a personal finance dashboard. A natural-language request went through the three layers and was released in version 1.0.0 of the product, with a human decision at every critical point, and the personalization agent adapted the interface per user and, in the worst case, rejected three malicious payloads that bypassed server-side validation. The validation also exposed real failures: half of the deliveries approved before the independent review was introduced were later reverted, and the cost estimates produced during planning were off by more than an order of magnitude. Each failure was absorbed by the architecture as a verifiable mechanism. The contributions of this work are the multidimensional analysis and its requirements, the architecture specification, the proof of concept and its evidence, and a design principle that emerged from the validation: guarantees that rely only on instructions to the model are not guarantees, and must be verified by code or decided by a person.

   \vspace{\onelineskip}

   \noindent
   \textbf{Keywords}: Agentic Artificial Intelligence. Software Engineering. Supervised Autonomy. Multi-Agent Systems. Adaptive Interfaces. AI Governance.
 \end{otherlanguage*}
\end{resumo}
